% Generated draft: requires mathematical and source review.
\documentclass[11pt]{article}
\usepackage[margin=25mm]{geometry}
\usepackage{fontspec}
\setmainfont{Latin Modern Roman}
\usepackage{amsmath,amssymb,mathtools,unicode-math}
\setmathfont{Latin Modern Math}
\usepackage{enumitem,needspace,xurl,hyperref}
\setlength{\parindent}{0pt}
\setlength{\parskip}{5pt}
\setlength{\emergencystretch}{3em}
\newcommand{\N}{\mathbb{N}}
\newcommand{\Z}{\mathbb{Z}}
\newcommand{\Q}{\mathbb{Q}}
\newcommand{\R}{\mathbb{R}}
\newcommand{\C}{\mathbb{C}}
\newcommand{\F}{\mathbb{F}}
\newcommand{\A}{\mathbb{A}}
\newcommand{\PP}{\mathbb{P}}
\newcommand{\E}{\mathbb{E}}
\newcommand{\eps}{\varepsilon}
\newcommand{\dd}{\,\mathrm{d}}
\begin{document}
\section*{UnsolvedMath ID: 701703}
\textbf{Source problem code:} catalogue entry
\subsection{Definitions and mathematical statement}
Let \(n\geq 1\) be an integer, and let \(G\leq \mathrm{GL}(n,\Z)\) be a finitely generated subgroup, where \(\mathrm{GL}(n,\Z)\) denotes the group of integer matrices with determinant \(\pm1\). For each integer \(b\geq1\), define
\[
e_b=\#\bigl\{(g_{ij})\in G: |g_{ij}|\leq b\text{ for every }1\leq i,j\leq n\bigr\}.
\]
This count is finite and positive, since there are finitely many integer matrices satisfying the entry bound and the identity matrix belongs to this set. Is it true that, for every such \(n\) and \(G\), there exists \(L\in\R\) such that
\[
\lim_{\substack{b\to\infty\\ b\in\Z,\ b\geq2}}\frac{\ln e_b}{\ln b}=L?
\]
Here \(\ln\) denotes the natural logarithm, and the bounds tend to infinity through integers.
\subsection{Short English statement}
For every finitely generated group of invertible integer matrices, count its matrices whose entries have absolute value at most \(b\). Must the logarithm of this count divided by \(\ln b\) converge as the integer bound \(b\) tends to infinity?
\subsection{Sources}
\begin{itemize}
\item Supplied local UnsolvedMath registry record: permanent ID \(701703\), exact source code \texttt{catalogue entry}, title \emph{Kourovka Notebook Problem 20.45}. Upstream revision: \texttt{d371ee917ac9647fbb06cd6fc529deec6f3dda82}.
\item \emph{The Kourovka Notebook (Unsolved Problems in Group Theory)}, Problem 20.45, proposed by S. Kohl in 2022. The supplied background identifies the 20th issue (2022) and the 21st issue edition (2026), \texttt{arXiv:1401.0300v48}.
\item The supplied collection description identifies Issues 1--21 (1965--2026), edited by E. I. Khukhro and V. D. Mazurov. The registry labels this problem open; that label is supplied source data and has not been independently verified.
\end{itemize}
\end{document}
